\section{Introduction}
\label{sec:intro}

Transformer layers evolve representations by adding learned attention and MLP
transformations to the current hidden state~\citep{vaswani2017attention,elhage2021mathematical}.
This additive form gives each transformation two geometrically different
roles: a parallel component that rescales the current state, and a
perpendicular component that redirects it by adding content in a new direction.
This view raises two questions about representation evolution. First, how much
of a learned transformation preserves the current representation direction,
and how much introduces direction-changing content? Second, if a parallel
component is non-negligible, does it affect model behavior, and can it be
modified or removed without disrupting the model? To answer these questions,
we apply the same geometry in two spaces. In \emph{residual space}, we
decompose the attention and MLP outputs that are added back to the hidden
state, measuring how each module evolves the representation. In \emph{value
space}, we analyze attention aggregation over projected values, providing an
observation space for token flow and context aggregation before the output
projection.

We first build layerwise profiles across pretrained models and find that
learned attention and MLP outputs contain substantial components parallel to
the current hidden state. This indicates that direction-preserving
transformations are not merely inherited from the additive identity path, even
though attention and MLP blocks are nominally responsible for mixing tokens
and transforming features.
However, layerwise profiles alone do not establish how these components affect
model behavior. We therefore use component-scaling experiments to test their
functional role. When we manipulate the part of the attention aggregation aligned
with the projected self-value, model behavior remains close to baseline over a
broad scaling range. In contrast, analogous residual-space attention scaling
is less stable, residual MLP-side scaling is more fragile, and perpendicular
scaling is much more destructive, consistent with the role of perpendicular
components in changing representation direction. Thus, \textit{parallel-component
robustness is space-dependent}: it depends on the computational space in
which that component is manipulated. This motivates a closer look at attention,
where token routing and value aggregation define a distinct space for such
manipulations.

To explain how token flow contributes to the parallel component, we apply the
same decomposition inside attention's value aggregation. This yields a
factorization of the attention-side parallel coefficient into diagonal
self-routing, value self-alignment, and cross-token terms. The diagonal term
has a simple interpretation: self-attention carries the token's own projected
value back into the aggregation, and this contribution becomes parallel when
that value is aligned with the current representation. The value-space view
therefore identifies the self-value direction as a natural coordinate for
manipulating this self-value-parallel contribution before the output
projection.

This distinction matters experimentally. Value-space parallel removal targets
the component aligned with the projected self-value while leaving the remaining
value aggregation intact. Residual-space analogues instead measure parallelness
after the output projection against a different reference direction. As a
result, nominally similar diagonal or parallel-removal edits can produce very
different zero-shot behavior, including in prior work on self-directed value
flow~\citep{zhai2025xsa}.

Finally, we evaluate whether this geometry extends beyond inference-time
interventions. In compression, common post-training quantization and
pruning settings introduce approximation error into layer
transformations~\citep{Frantar2022GPTQAP,Lin2023AWQAW,Sun2023ASA}.
We find that stronger compressed models have lower perpendicular error in
their layer transformations, while parallel error is less discriminative. This
yields a label-free diagnostic signal for compression: preserving
perpendicular update geometry is more predictive of preserved behavior than
preserving the parallel component. The relative robustness of value-space
parallel manipulation also motivates a training-time test. In retained
pretraining runs, suppressing attention-side self-value-parallel components
changes the optimization trajectory and can improve average downstream scores.
This is consistent with the view that self-value-parallel suppression can
shift training pressure toward more direction-changing updates.
Together, these results suggest that parallel/perpendicular geometry is not
just a descriptive decomposition, but \textit{an organizing principle} for safe
editing, compression diagnostics, and training-time intervention.

\paragraph{Contributions.}
\begin{itemize}
  \item This work reveals that substantial direction-preserving parallel
        components are widespread in learned attention and MLP transformations.
        We further propose a two-space geometric decomposition in residual and
        value spaces to measure their role in representation evolution and
        localize their attention-side source.
  \item In inference-time interventions, we identify a clear space dependence:
        value-space manipulation remains closest to baseline, whereas
        residual-space scaling is more fragile and perpendicular scaling is
        substantially more disruptive.
  \item The attention-side factorization decomposes the parallel contribution
        into diagonal self-routing, value self-alignment, and cross-token
        aggregation. This separates the self-value component from the
        remaining context aggregation and explains why value-space and
        residual-space edits lead to different zero-shot behavior.
  \item Beyond inference-time editing, the same geometry provides practical
        guidance for compression and training: lower perpendicular
        transformation error tracks stronger compressed models, and
        self-value-parallel suppression in retained pretraining runs shifts
        optimization and can improve downstream performance.
\end{itemize}
